\documentclass[10pt,a4paper]{amsart}
\usepackage[margin=19mm]{geometry}
\usepackage{amsmath,amssymb,mathtools}
\usepackage[scr=rsfs,cal=euler]{mathalfa}
\usepackage{xfrac}
\usepackage[unicode,hidelinks,bookmarks=false]{hyperref}
\newcommand{\ie}{i.e.\ }
\newcommand{\cf}{cf.\ }
\newcommand{\resp}{respectively\ }
\newcommand{\Subsection}{\subsection}
\providecommand{\nobreakdash}{\nobreak\hbox{-}}
\setlength{\emergencystretch}{2em}
\multlinegap0pt
\usepackage{csquotes}
\usepackage{amssymb}
\usepackage{bm}
\usepackage{tcolorbox}
\usepackage{xcolor}
\newtheorem{proposition}{Proposition}
\title{A Group with Exactly One Noncommutator}
\author{Omar Hatem, Daoud Siniora}
\date{Source: arXiv:2511.00541v2 (CC BY 4.0)}
\begin{document}
\maketitle

\noindent\emph{Source and licence.} A Group with Exactly One Noncommutator. Version arXiv:2511.00541v2 (CC BY 4.0), distributed by arXiv under the Creative Commons Attribution 4.0 International licence, http://creativecommons.org/licenses/by/4.0/, as recorded in the arXiv licence metadata for this exact version and shown on the abstract page https://arxiv.org/abs/2511.00541v2. Full licence notice: \texttt{licenses/arxiv-2511.00541-CC-BY-4.0.txt}.\par\smallskip

\noindent\textit{Editorial scope.} Counted unit: the article's proposition four\_props (source0.tex lines 94--103). The article's complete original proof of it is delivered with it (source0.tex lines 104--111, one \texttt{proof} environment of 8 lines). No mathematics is written, repaired or re-derived: the statement and the proof are the article's own lines, in the article's own notation, and no retained line is substituted or re-flowed.  No further source-local context is needed: every cross-reference of every retained line resolves inside the two delivered spans.  Nothing was omitted from any retained span, so the package carries no omission record.  No retained line contains a citation, so the wrapper carries no bibliography and no bibliography entry.  No retained line contains a figure environment, an image-inclusion command or drawing code, so no image is delivered and the package declares no asset dependency.  Editorial formatting only, in addition: an AMS \texttt{amsart} preamble that restates the article's own declarations and macros used by the retained lines, namely the environments \texttt{proposition} on separate counters, together with the standard packages those lines need.  The retained environments print in this document as Proposition 1 in order of appearance, whereas the article prints the same items with the numbering of its own full document.  The complete native source of the article as distributed in the arXiv e-print for this exact version is bundled unchanged as \texttt{sources/188256/source0.tex}.
\begin{proposition}
Let $G$ be a finite group with $|G|\geq 3$ that has exactly one noncommutator element $u\in G$. Then the following hold.  
\begin{itemize}
            \item The order of $u$ is $2$.
            \item The element $u$ belongs to the center $Z(G)$ of $G$.
            \item $G$ is a perfect group.
            \item The element $u$ is the product of two commutators. Consequently, the commutator length of $G$ is $2$.
\end{itemize}
\label{four_props}
\end{proposition}

\begin{proof} Let $u$ be the unique noncommutator element of $G$.
\begin{itemize}
    \item The inverse of a commutator is a commutator, and so the inverse of a noncommutator is a noncommutator. Since $u$ is unique, it equals its inverse. Thus, $u^2 = 1$.
    \item Since conjugation by $g\in G$ is an automorphism of $G$, it follows that $gug^{-1}$ must be a noncommutator as well. Since $u$ is the only noncommutator it must be $gug^{-1}=u$ for any $g\in G$.  Therefore, for any $g\in G$, we have $gu=ug$, which means $u \in Z(G)$.
    \item Let $S$ be the set of all commutators of $G$, so $|S|+1=|G|$. Since $G'=\langle S \rangle\subseteq G$ we get that $|G'|\geq |S| = |G|-1\geq 2$ and either $G'=S$ or $G'=G$. But $|G'|$ divides $|G|$, and thus $|G|=|G'|$ implying that $G=G'$ showing that $G$ is perfect.
    \item Since $|G| > 2$, we can choose some nontrivial commutator $g\in G$. Then $u=(ug)g^{-1}$. Since $g$ is a commutator, we know $g^{-1}$ is a commutator as well. If $ug$ were a noncommutator, then $ug=u$ because $u$ is the only noncommutator, and so $g=1$ contradicting that $g$ is nontrivial. So $ug$ is a commutator too. Therefore, $u$ is a product of two commutators, and thus the commutator length of $G$ is $2$.
\end{itemize}
\end{proof}

\end{document}
